\begin{frame}[fragile]{Dyna-Q 코드: 경험 $\to$ 학습 $\to$ 모델 저장 $\to$ 계획}
  \begin{lstlisting}[basicstyle=\ttfamily\scriptsize]
def dyna_q(env_step, n_states, n_actions, n_episodes,
           planning_steps, alpha, gamma, epsilon, start_state):
    Q = [[0.0] * n_actions for _ in range(n_states)]
    model = {}  # (s, a) -> (r, s')  관찰한 것을 그대로 저장
    for _ in range(n_episodes):
        s = start_state
        for _ in range(100):
            a = random.randrange(n_actions) if random.random() < epsilon \
                else max(range(n_actions), key=lambda x: Q[s][x])
            ns, r, done = env_step(s, a)
            Q[s][a] += alpha * (r + gamma * max(Q[ns]) - Q[s][a])  # 1) 실제 학습
            model[(s, a)] = (r, ns)                                 # 2) 모델 갱신
            for _ in range(planning_steps):                         # 3) 계획
                (ps, pa), (pr, pns) = random.choice(list(model.items()))
                Q[ps][pa] += alpha * (pr + gamma * max(Q[pns]) - Q[ps][pa])
            s = ns
            if done:
                break
    return Q
  \end{lstlisting}
  {\footnotesize 결정론적 환경 가정(같은 $(s,a)$은 항상 같은
  $(r,s')$을 냄). \texttt{planning\_steps}가 클수록 실제 스텝 1회당
  경험을 모델을 통해 ``몇 번이고 복습''.}
\end{frame}
